% Lösungen Einheit 1 von 5 – Tangentengleichung im Punkt (zusammenbau v0.1)
\einheitenkopf{Einheit 1 von 5 · Tangentengleichung im Punkt}

\erg{9}{a) ein Punkt – ableiten und die Stelle einsetzen. \quad b) $f(2) = 4$, $f'(x) = 3x^2 - 2$, $m = f'(2) = 10$; $n = 4 - (10) \cdot (2) = -16$; $t(x) = 10x - 16$ \quad c) $f(-2) = 8 + 8 - 6 - 1 = 9$ – $P$ liegt auf dem Graphen; $f'(-2) = -12 - 8 + 3 = -17$; $n = 9 - 34 = -25$; $t(x) = -17x - 25$ \quad d) waagerecht je $3$ m (halbe Breite): links $m = \frac{2}{3}$, rechts $m = -\frac{2}{3}$; Länge $= \sqrt{3^2 + 2^2} = \sqrt{13} \approx 3{,}61$ m \quad e) $f(0) = 2 - 3 = -1$; $f'(x) = 2e^x$, $m = f'(0) = 2e^0 = 2$ (nicht $0$); $t(x) = 2x - 1$ \quad f) $f'(x) = 3x^2 - 6x = 9 \Leftrightarrow x^2 - 2x - 3 = 0$; $x_1 = 3$, $x_2 = -1$; Punkte $(3 | 1)$ und $(-1 | -3)$ \quad g) $(-1 + \sqrt{6})^2 = 7 - 2\sqrt{6}$; $f'(-1 + \sqrt{6}) = (7 - 2\sqrt{6} - 2 + 2\sqrt{6} - 5) \cdot e^{-1 + \sqrt{6}} = 0$; $f(-1 + \sqrt{6}) = (2 - 2\sqrt{6}) \cdot e^{\sqrt{6} - 1} \approx -12{,}35$ \quad h) $f'(x) = 3x^2 - 3 = 9$ liefert $x = -2$; Berührpunkt $(-2 | -2)$, Tangente $y = 9x + 16$, etwa durch $(-2 | -2)$ und $(-1{,}5 | 2{,}5)$ \quad i) zwei Gitterpunkte, etwa $(2 | 0)$ und $(4 | 4)$: $m = \frac{4 - 0}{4 - 2} = 2$; $t(x) = 2x - 4$ \quad j) $f(1) = -2$, $f(-2) = 4$; Steigung von $g$: $\frac{4 - (-2)}{-2 - 1} = -2$; $f'(x) = 4x^3 - 6x$, $f'(1) = -2$ – gleiche Steigung, und $g$ geht durch den Berührpunkt: $g$ ist Tangente. \quad k) $m = f_a'(u) = 3au^2$; $n = f_a(u) - m \cdot u = au^3 - 3au^3 = -2au^3 = -2f_a(u)$ – allgemein für jedes $u$, nicht nur an einem Beispiel.}

9h)

\begin{ksys}[xmin=-3,xmax=3,ymin=-5,ymax=5,klein] \funktion{\x^3-3*\x}{f} \gerade{9}{16}{t} \end{ksys}

\erg{10}{a) $n(c) = f(c) - c \cdot f'(c) = \frac{c^2 + 16 - c^2}{\sqrt{c^2 + 16}} = \frac{16}{\sqrt{c^2 + 16}}$; $n(c) > 2 \Leftrightarrow \sqrt{c^2 + 16} < 8 \Leftrightarrow c^2 < 48$; Grenzen $c \approx -6{,}93$ und $c \approx 6{,}93$, also $-6{,}93 < c < 6{,}93$}
\erg{11}{a) $f_k'(x) = 3k^2x^2 - 6kx + 4$, $f_k''(x) = 6k^2x - 6k = 0 \Leftrightarrow x = \frac{1}{k}$; $f_k'(\frac{1}{k}) = 3 - 6 + 4 = 1$ – unabhängig von $k$, also sind alle Wendetangenten parallel.}
\erg{12}{a) Ben nimmt den Funktionswert als Steigung. Richtig: $m = f'(1) = 2 + 3 = 5$; aus $4 = 5 + n$ folgt $n = -1$; $t(x) = 5x - 1$. \quad b) Die Steigung des Graphen an einer Stelle ist als Steigung seiner Tangente dort festgelegt – beides ist der Ableitungswert $f'(x_0)$.}
